\clearpage
\section{The universal mixed-weight jet filter}\label{sec:jet}
\subsection{One polynomial chooses the cancelling coefficients}
Let $m\ge0$ be an integer. Define
\[
 P_m(X)=X\prod_{j=1}^m(jX-1)=\sum_{s=1}^{m+1}c_{m,s}X^s,
\]
and the formal series
\[
 \Phi_m(t)=\frac1{(m+1)!}\sum_{s=1}^{m+1}c_{m,s}\Li_s(t).
\]
Here $\Li_s(t)=\sum_{k\ge1}t^k/k^s$ is first interpreted in $\Q[[t]]$. It also defines the usual analytic polylogarithm on $|t|<1$~\cite{DLMFPolylog}. The label $s$ is its weight. The contact order will follow from its coefficients, rather than from that label.

\begin{theorem}[Exact jet cancellation]\label{thm:jet-contact}
\[
 \Phi_m(t)=(-1)^m\sum_{k\ge1}\binom{k}{m+1}\frac{t^k}{k^{m+2}}.
\]
Thus
\[
 \ord_0\Phi_m=m+1,\qquad
 [t^{m+1}]\Phi_m=\frac{(-1)^m}{(m+1)^{m+2}}.
\]
\end{theorem}
\begin{proof}
The coefficient of $t^k$ is $P_m(1/k)/(m+1)!$. Direct substitution gives
\[
 P_m(1/k)=\frac1k\prod_{j=1}^m\left(\frac jk-1\right)
 =(-1)^m\frac{(k-1)\cdots(k-m)}{k^{m+1}}.
\]
Divide by $(m+1)!$ and multiply numerator and denominator by $k$ to get the asserted binomial coefficient. It vanishes for $1\le k\le m$ and first becomes one at $k=m+1$. The coefficient there is nonzero.
\end{proof}
For $m=0$, this is $\Phi_0=\Li_1$. For $m=1$,
\[
 \Phi_1=(\Li_2-\Li_1)/2=-t^2/8-t^3/9-\cdots.
\]
The sign is part of the exact formula. Changing the weight alone would not produce the cancellation of the linear coefficient.

\subsection{Small coefficient norm has a limited meaning}
The coefficients of $P_m$ alternate in sign. Its absolute-coefficient polynomial evaluated at one gives
\[
 \sum_s|c_{m,s}|=\prod_{j=1}^m(j+1)=(m+1)!.
\]
Also $c_{m,1}=(-1)^m$, so the integer coefficient vector is primitive. After the displayed normalization, its coefficient $\ell^1$ norm is one.

These are exact statements about a vector of rational coefficients. The number of participating weights grows with $m$. A norm of one for the coefficient vector does not bound the denominator, covolume or arithmetic height of a period realization. No such realization is built into the definition of $\Phi_m$.

\subsection{The depleted local series}
For a prime $p$, delete the terms whose index is divisible by $p$:
\[
 \Phi_{m,p}(t)=(-1)^m\sum_{p\nmid k}\binom{k}{m+1}\frac{t^k}{k^{m+2}}\in\Z_p[[t]].
\]
It converges when $v_p(t)>0$.
\begin{theorem}[Exact local jet valuation]
If $v_p(t)=e\ge1$, then
\[
 v_p(\Phi_{m,p}(t))=
 \begin{cases}
 (m+1)e,&p\nmid m+1,\\
 (m+2)e,&p\mid m+1.
 \end{cases}
\]
\end{theorem}
\begin{proof}
If $p\nmid m+1$, the first nonzero coefficient remains the unit $(-1)^m/(m+1)^{m+2}$. Its term has valuation $(m+1)e$. Every later coefficient is $p$-integral, so every later term has strictly larger valuation.

If $p\mid m+1$, that first index is deleted. The next index is $m+2$, which is a $p$-unit. Its coefficient is $(-1)^m/(m+2)^{m+1}$, also a unit. All later terms again have strictly larger valuation. The ultrametric inequality with a unique least-valuation term proves both equalities.
\end{proof}

\subsection{An exact differential identity}
Let $D=t\,\dd/\dd t$ be the Euler operator. Termwise formal differentiation is valid. Applying $D^{m+2}$ multiplies the coefficient at $t^k$ by $k^{m+2}$, so
\[
 D^{m+2}\Phi_m=(-1)^m\sum_{k\ge1}\binom{k}{m+1}t^k
 =(-1)^m\frac{t^{m+1}}{(1-t)^{m+2}}.
\]
The final equality is the elementary generating function for the binomial coefficients.

This universal formal identity and the prime-dependent depleted valuations can be studied simultaneously. They do not identify the different $\Phi_{m,p}$ as local values of one algebraic number. A product formula cannot be applied merely because these functions arose from one coefficient pattern.
